\section{Fazit und Ausblick}
\label{sec:fazit}

\subsection{Zusammenfassung}

Im Rahmen dieses Projekts wurde ein Feedforward-Netz vollständig von Grund auf in Python implementiert -- ohne Rückgriff auf fertige Deep-Learning-Bibliotheken. Alle Komponenten -- Forward Pass, Loss-Berechnung, Backpropagation und Gewichtsupdate -- wurden selbst geschrieben und sind in diesem Dokument dokumentiert.

Die drei Ziele aus der Einleitung wurden erreicht:

\begin{enumerate}
    \item \textbf{Theoretisches Verständnis:} Die mathematischen Grundlagen sind in \autoref{sec:theorie} ausführlich dargestellt.
    \item \textbf{Implementierung:} Das Netz unterstützt beliebig viele Hidden Layers sowie Gewichte und Biases (siehe \autoref{sec:implementierung}).
    \item \textbf{Evaluation:} Das Netz wurde auf einem Datensatz trainiert und die Ergebnisse in \autoref{sec:ergebnisse} dokumentiert.
\end{enumerate}


\subsection{Persönliches Fazit}

\begin{hinweisbox}[Platzhalter]
Kurzer Absatz jedes Teammitglieds: Was habt ihr persönlich gelernt? Was würdet ihr beim nächsten Mal anders machen?
\end{hinweisbox}
